Аллокаторы (\textit{allocator} "--- распределитель памяти) используются с целью изоляции разработчика алгоритмов и контейнеров, нуждающихся в выделении памяти, от низкоуровневых деталей физической организации памяти. Аллокаторы предоставляют стандартные способы выделения и освобождения памяти, а также стандартные имена типов, используемые для указателей и ссылок. Принципиально аллокатор является абстракцией, которая позволяет разработчикам контролировать стратегии распределения памяти без изменения кода контейнеров.

Бьёрн Страуструп, создатель языка C++, противопоставляет итераторы и аллокаторы, говоря о том, что итераторы, поддерживающие абстрактную модель данных в виде последовательности объектов, "--- это понятие, с которым должен быть знаком каждый программист, тогда как аллокаторы, обеспечивающие отображение низкоуровневой последовательности байтов в высокоуровневую объектную модель, "--- механизм, сведения о работе которого в большинстве случаев не нужны программисту, ведь стандартная библиотека C++ предоставляет стандартный аллокатор, удовлетворяющий нужды большинства пользователей \cite{straustrup}.

Тем не менее, у пользователей также есть возможность разрабатывать собственные аллокаторы, чтобы реализовать альтернативные схемы работы с памятью, такие как использование разделяемой памяти, памяти с автоматической сборкой мусора или памяти на базе заранее подготовленного пула объектов. В некоторых задачах пользовательский аллокатор может значительно выиграть в скорости по сравнению со стандартным. Именно поэтому \textbf{целью моей работы} является проведение сравнительного анализа различных алгоритмов аллокации памяти и реализация собственного аллокатора. 

\textbf{Задачи работы} включают:
\begin{itemize}
	\item изучение реализации стандартного аллокатора в языке C++;
	\item проведение обзора и классификации существующих алгоритмов аллокации памяти и реализацию некоторых из этих алгоритмов на языке C++;
	\item создание и реализацию собственного алгоритма аллокации памяти на базе существующих с применением некоторых оптимизаций;
	\item выполнение сравнительного анализа существующих алгоритмов, а также собственного алгоритма, с целью сравнения их быстродействия.
\end{itemize}
